\section{Limitations}
\label{sec:limitations}

Only two health datasets were available. They differ substantially in size and outcome context, which makes the contrast informative, but two datasets cannot identify a general relationship between intrinsic overlap and SepAware benefit. The current paper should therefore be read as an initial multidataset validation with a controlled mechanism experiment. Additional health datasets with independently observed overlap, prevalence, sample size, and minority heterogeneity are required.

The strongest Vigitel factorial finding comes from one 5000-record sample, one cross-validation partition, one CTGAN configuration, and shared within-fold candidate pools. The strengthened experiment uses three repeated five-fold partitions and 30 training epochs for computational feasibility. Fold-specific generator seeds avoid reusing a single fitted generator, and repeated shuffles measure some partition variability, but all repetitions still reuse the same two source samples. More generator seeds, higher-capacity sensitivity runs, and independent datasets are necessary before formal confirmatory claims.

The fixed lambda grid was exploratory because performance and configuration selection used the same folds. The strengthened analysis fixes the separability weight in advance, but it does not estimate a clinically optimal trade-off between discrimination and minority diversity. The score also depends on a random-forest confidence term; additional boundary models and label-permutation controls are needed to test classifier-specific shortcuts.

The TVAE condition is not directly symmetric with CTGAN. SDV 1.17 does not implement native TVAE conditional sampling, so TVAE was fitted to minority-class features. This is a reproducible conditional-generation strategy, but its low diversity shows that post-selection cannot compensate for a collapsed or highly concentrated candidate pool. A diffusion model and a generator with native conditional support should be added before making a broad generator-agnostic claim.

Minority coverage and diversity are geometric proxies rather than clinically validated phenotype measures. The COVID-19 death class may contain several pathways that are not represented by unsupervised distances. Blinded clinician review and prespecified phenotype or subgroup definitions are required to determine whether removed variation is clinically meaningful. Calibration estimates are also unstable in the small COVID-19 folds.

The privacy analysis is limited to exact matches, nearest-record distances, and a membership proxy. It does not include attribute inference, shadow-model membership attacks, linkage attacks, or a formal differential-privacy guarantee. No synthetic data release should be justified from the present privacy results.

Finally, external or temporal validation is absent. Vigitel spans multiple survey years, but the current evaluation is not a prospective temporal split. The COVID-19 cohort is small and institution-specific. At least one temporal or cross-site evaluation is needed to establish robustness under distribution shift.

These limitations define the next additions to the manuscript: 2--4 further health datasets, repeated partition and generator seeds, a native conditional third generator, expert-informed minority coverage, stronger privacy attacks, and external or temporal validation. If SepAware fails under some of these conditions, those failures should be reported as operational guidance about when separability control should not be used.
